\documentclass[12pt,a4paper]{article}
% Краткое сообщение Ксении по итогам ревью: в задании 1 ошибка возникла при возведении в квадрат отклонения для среднего значения; в задании 2 численный результат верный, но не использован прямо требуемый условием способ; в задании 6 точное значение дисперсии найдено верно, но десятичная запись дроби указана неверно.
\usepackage[a4paper,margin=2cm]{geometry}
\usepackage[T2A]{fontenc}
\usepackage[utf8]{inputenc}
\usepackage[russian]{babel}
\usepackage{amsmath,amssymb,mathtools}
\usepackage{array,booktabs,longtable}
\usepackage{xcolor}
\usepackage{hyperref}
\usepackage{tikz}

\hypersetup{
  colorlinks=true,
  linkcolor=blue!55!black,
  urlcolor=blue!55!black
}
\setlength{\parindent}{0pt}
\setlength{\parskip}{0.6em}
\renewcommand{\arraystretch}{1.25}

\begin{document}

\begin{titlepage}
\thispagestyle{empty}
\begin{center}
\vspace*{0.27\textheight}
{\LARGE\bfseries Ревью домашней работы\par}
\vspace{1.5cm}
{\large Автор: Лёвин Артём Александрович\par}
\vspace{0.5cm}
{\large 5 октября 2026 года\par}
\vfill
\begin{tikzpicture}
\draw[blue!30,line width=0.9pt]
  (0,0) .. controls (1.2,0.35) and (2.4,-0.35) .. (3.6,0);
\end{tikzpicture}
\end{center}
\end{titlepage}

\newpage
\section*{Сводная таблица}
\small
\begin{longtable}{p{0.055\linewidth}p{0.17\linewidth}p{0.265\linewidth}p{0.15\linewidth}p{0.15\linewidth}}
\toprule
\textbf{№} & \textbf{Тема} & \textbf{Суть ошибки} & \textbf{Ваш ответ} & \textbf{Правильный ответ}\\
\midrule
\endfirsthead
\toprule
\textbf{№} & \textbf{Тема} & \textbf{Суть ошибки} & \textbf{Ваш ответ} & \textbf{Правильный ответ}\\
\midrule
\endhead
\hyperref[task:1]{1} & Математическое ожидание и дисперсия & Для значения $X=2$ квадрат отклонения от среднего принят равным $4$ вместо $0$, поэтому дисперсия получилась завышенной. & $M(X)=2$, $D(X)=4$ & $M(X)=2$, $D(X)=2$\\
\midrule
\hyperref[task:2]{2} & Дисперсия по формуле $M(X^2)-(M(X))^2$ & Численный результат верен, но не выполнено прямое требование условия использовать указанную формулу. & $M(X)=0{,}4$, $D(X)=0{,}84$ & $M(X)=0{,}4$, $D(X)=0{,}84$\\
\midrule
\hyperref[task:6]{6} & Несмещённая оценка дисперсии & После верного точного значения $s^2=\frac53$ записано неверное равенство $\frac53=1{,}6$. & $s^2=1{,}6$ & $s^2=\frac53\approx1{,}67$\\
\bottomrule
\end{longtable}
\normalsize

\newpage
\thispagestyle{empty}
\vspace*{\fill}
\begin{center}
{\LARGE\bfseries Разбор ошибок}
\end{center}
\vspace*{\fill}

\newpage
\phantomsection\label{task:1}
\section*{Задание 1}

\textbf{Текст задания.}
Случайная величина $X$ принимает значения $0$, $2$ и $4$ с вероятностями $0{,}25$, $0{,}5$ и $0{,}25$. Найдите $M(X)$ и $D(X)$.

\textbf{Ваше решение.}
Математическое ожидание найдено верно:
\[
M(X)=0\cdot0{,}25+2\cdot0{,}5+4\cdot0{,}25=2.
\]
Далее были выписаны отклонения от математического ожидания. Для крайних значений получены $-2$ и $2$, а для среднего значения $2$ отклонение равно $0$. После этого при возведении отклонений в квадрат в записи появились три значения $4$, и дисперсия была посчитана как
\[
D(X)=4\cdot0{,}25+4\cdot0{,}5+4\cdot0{,}25=4.
\]

\textbf{Ваш ответ.}
$M(X)=2$, $D(X)=4$.

\textcolor{red}{\textbf{Ошибка:}} неверно найден квадрат отклонения для значения $X=2$.

\textbf{Где именно возникает ошибка.}
До шага
\[
2-M(X)=2-2=0
\]
решение идёт верно. Первый неверный переход появляется при возведении этого отклонения в квадрат: вместо
\[
0^2=0
\]
в расчёт попадает число $4$.

\textbf{Почему это неверно.}
Дисперсия дискретной случайной величины по определению через отклонения строится из квадратов расстояний каждого возможного значения до математического ожидания. Здесь математическое ожидание равно $2$. Значение $X=2$ совпадает со средним, поэтому его расстояние до среднего равно нулю. Квадрат нуля также равен нулю. Вероятность $0{,}5$ относится именно к этому нулевому отклонению и потому не должна давать вклад $4\cdot0{,}5$ в дисперсию.

\textcolor{blue!60!black}{\textbf{Ключевая идея:}} если одно из значений случайной величины совпадает с математическим ожиданием, его вклад в дисперсию равен нулю независимо от его вероятности.

\textbf{Исправление от последнего правильного шага.}
Отклонения равны
\[
0-2=-2,\qquad 2-2=0,\qquad 4-2=2.
\]
Их квадраты:
\[
(-2)^2=4,\qquad 0^2=0,\qquad 2^2=4.
\]
Поэтому
\[
D(X)=4\cdot0{,}25+0\cdot0{,}5+4\cdot0{,}25=1+0+1=2.
\]

\textbf{Полное правильное решение.}
Сначала находим математическое ожидание:
\[
M(X)=0\cdot0{,}25+2\cdot0{,}5+4\cdot0{,}25=0+1+1=2.
\]
Затем вычисляем квадрат отклонения каждого значения от найденного среднего:
\[
(0-2)^2=4,
\]
\[
(2-2)^2=0,
\]
\[
(4-2)^2=4.
\]
Умножаем каждый квадрат отклонения на соответствующую вероятность и складываем:
\[
D(X)=4\cdot0{,}25+0\cdot0{,}5+4\cdot0{,}25=2.
\]
Следовательно,
\[
M(X)=2,\qquad D(X)=2.
\]

\textbf{Проверка.}
Проверим дисперсию второй формулой. Сначала
\[
M(X^2)=0^2\cdot0{,}25+2^2\cdot0{,}5+4^2\cdot0{,}25
=0+2+4=6.
\]
Тогда
\[
D(X)=M(X^2)-(M(X))^2=6-2^2=6-4=2.
\]
Получено то же значение.

\textcolor{teal!60!black}{\textbf{Итог:}} математическое ожидание найдено верно, а дисперсия равна $2$, потому что для значения $X=2$ квадрат отклонения равен нулю.

\textbf{Задача на закрепление.}
Случайная величина $Y$ принимает значения $1,3,5$ с вероятностями $0{,}25;\ 0{,}5;\ 0{,}25$. Найдите $M(Y)$ и $D(Y)$.

\newpage
\phantomsection\label{task:2}
\section*{Задание 2}

\textbf{Текст задания.}
Случайная величина $X$ принимает значения $-1$ и $1$ с вероятностями $0{,}3$ и $0{,}7$. Найдите $M(X)$ и $D(X)$, используя формулу
\[
D(X)=M(X^2)-(M(X))^2.
\]

\textbf{Ваше решение.}
Математическое ожидание вычислено верно:
\[
M(X)=-1\cdot0{,}3+1\cdot0{,}7=0{,}4.
\]
После этого дисперсия найдена через квадраты отклонений от математического ожидания:
\[
(-1-0{,}4)^2=1{,}96,\qquad (1-0{,}4)^2=0{,}36,
\]
\[
D(X)=1{,}96\cdot0{,}3+0{,}36\cdot0{,}7=0{,}84.
\]

\textbf{Ваш ответ.}
$M(X)=0{,}4$, $D(X)=0{,}84$.

\textcolor{red}{\textbf{Ошибка:}} не выполнено требование использовать указанную в условии формулу.

\textbf{Где именно возникает ошибка.}
Численные вычисления не содержат ошибки: и математическое ожидание, и дисперсия получены верно. Несоответствие возникает на этапе выбора способа вычисления дисперсии. В условии прямо сказано использовать формулу через $M(X^2)$, а в решении применён другой, хотя и равносильный, способ --- через средний квадрат отклонения.

\textbf{Почему это неверно.}
Если в условии указан конкретный метод, то часть задания состоит не только в получении правильного числа, но и в демонстрации владения именно этим методом. Поэтому правильный численный ответ не снимает требование вычислить $M(X^2)$ и затем вычесть квадрат $M(X)$.

\textcolor{blue!60!black}{\textbf{Ключевая идея:}} здесь особенно удобно заметить, что при $X=-1$ и $X=1$ квадрат случайной величины всегда равен $1$, поэтому $M(X^2)$ находится сразу.

\textbf{Исправление от последнего правильного шага.}
Уже найдено
\[
M(X)=0{,}4.
\]
Теперь используем именно требуемую формулу. Поскольку
\[
(-1)^2=1,\qquad 1^2=1,
\]
получаем
\[
M(X^2)=1\cdot0{,}3+1\cdot0{,}7=1.
\]
Тогда
\[
D(X)=M(X^2)-(M(X))^2=1-(0{,}4)^2=1-0{,}16=0{,}84.
\]

\textbf{Полное правильное решение.}
Находим математическое ожидание:
\[
M(X)=-1\cdot0{,}3+1\cdot0{,}7=-0{,}3+0{,}7=0{,}4.
\]
Далее находим математическое ожидание квадрата случайной величины:
\[
M(X^2)=(-1)^2\cdot0{,}3+1^2\cdot0{,}7
=1\cdot0{,}3+1\cdot0{,}7=1.
\]
Подставляем в указанную в условии формулу:
\[
D(X)=M(X^2)-(M(X))^2=1-(0{,}4)^2=0{,}84.
\]
Следовательно,
\[
M(X)=0{,}4,\qquad D(X)=0{,}84.
\]

\textbf{Проверка.}
Через квадраты отклонений получаем
\[
(-1-0{,}4)^2\cdot0{,}3+(1-0{,}4)^2\cdot0{,}7
=1{,}96\cdot0{,}3+0{,}36\cdot0{,}7
=0{,}84.
\]
Результат совпадает.

\textcolor{teal!60!black}{\textbf{Итог:}} численный ответ $0{,}84$ верен, но для полного выполнения задания нужно было получить его через $M(X^2)-(M(X))^2$.

\textbf{Задача на закрепление.}
Случайная величина $Y$ принимает значения $-2$ и $2$ с вероятностями $0{,}25$ и $0{,}75$. Найдите $M(Y)$ и $D(Y)$, используя формулу $D(Y)=M(Y^2)-(M(Y))^2$.

\newpage
\phantomsection\label{task:6}
\section*{Задание 6}

\textbf{Текст задания.}
Для выборки $1,2,3,4$ найдите несмещённую оценку дисперсии $s^2$.

\textbf{Ваше решение.}
Среднее найдено верно:
\[
\overline{x}=\frac{1+2+3+4}{4}=2{,}5.
\]
Сначала вычислена дисперсия с делением на $4$:
\[
D_n=\frac{(1-2{,}5)^2+(2-2{,}5)^2+(3-2{,}5)^2+(4-2{,}5)^2}{4}=\frac54.
\]
Затем выполнена поправка на несмещённость:
\[
s^2=\frac{4}{3}\cdot\frac54=\frac53.
\]
После этого записано
\[
\frac53=1{,}6,
\]
а затем дополнительно вычислялось приближённое значение стандартного отклонения.

\textbf{Ваш ответ.}
$s^2=1{,}6$.

\textcolor{red}{\textbf{Ошибка:}} неверно переведена дробь $\frac53$ в десятичную запись.

\textbf{Где именно возникает ошибка.}
Последний полностью правильный шаг:
\[
s^2=\frac53.
\]
Первый неверный шаг:
\[
\frac53=1{,}6.
\]

\textbf{Почему это неверно.}
Деление $5$ на $3$ даёт бесконечную периодическую десятичную дробь:
\[
\frac53=1{,}666\ldots
\]
Поэтому число $1{,}6$ не равно $\frac53$. Если оставлять точный ответ, нужно написать $\frac53$. Если требуется десятичное приближение до сотых, получаем $1{,}67$; до десятых --- $1{,}7$.

Дополнительное вычисление $s$ в этом задании не требуется, потому что условие просит именно несмещённую оценку дисперсии $s^2$.

\textcolor{blue!60!black}{\textbf{Ключевая идея:}} точная дробь $\frac53$ уже является полноценным окончательным ответом; при переходе к десятичной записи знак равенства можно сохранять только для точного десятичного представления, а для округления нужен знак приближённого равенства.

\textbf{Исправление от последнего правильного шага.}
Оставляем полученный верный результат:
\[
s^2=\frac53.
\]
При необходимости даём приближение:
\[
s^2=\frac53\approx1{,}67.
\]

\textbf{Полное правильное решение.}
Объём выборки равен $n=4$. Сначала находим среднее:
\[
\overline{x}=\frac{1+2+3+4}{4}=2{,}5.
\]
Сумма квадратов отклонений от среднего равна
\[
(1-2{,}5)^2+(2-2{,}5)^2+(3-2{,}5)^2+(4-2{,}5)^2
\]
\[
=2{,}25+0{,}25+0{,}25+2{,}25=5.
\]
Для несмещённой оценки дисперсии делим эту сумму на $n-1=3$:
\[
s^2=\frac{5}{3}.
\]
Следовательно,
\[
s^2=\frac53\approx1{,}67.
\]

\textbf{Проверка.}
Через промежуточную выборочную дисперсию с делением на $n$:
\[
D_n=\frac54.
\]
Поправка на несмещённость даёт
\[
\frac{n}{n-1}D_n=\frac43\cdot\frac54=\frac53,
\]
что совпадает с прямым вычислением через деление суммы квадратов отклонений на $3$.

\textcolor{teal!60!black}{\textbf{Итог:}} вычислительная схема была выстроена правильно до самого конца; точный ответ равен $\frac53$, а его десятичное приближение --- $1{,}67$.

\textbf{Задача на закрепление.}
Для выборки $2,3,5,6$ найдите несмещённую оценку дисперсии $s^2$.

\vfill
\begin{center}
\small Лёвин Артём Александрович эксклюзивно для Ксении Клыковой
\end{center}

\end{document}
